\chapter{Burndown charts des sprints}
\label{anx:burndowns}

Les burndown charts ci-dessous représentent l'évolution du nombre
de story points restants par jour, comparée à la trajectoire idéale.
Les zones rouges indiquent un retard ponctuel par rapport au plan
initial, généralement résorbé en fin de sprint.

\begin{figure}[H]\centering
\includegraphics[width=0.85\linewidth]{burndown-s1}
\caption{Burndown du sprint 1 --- Fondations RAG.}\label{fig:burndown-s1}
\end{figure}

\begin{figure}[H]\centering
\includegraphics[width=0.85\linewidth]{burndown-s2}
\caption{Burndown du sprint 2 --- Priorisation v1.}\label{fig:burndown-s2}
\end{figure}

\begin{figure}[H]\centering
\includegraphics[width=0.85\linewidth]{burndown-s3}
\caption{Burndown du sprint 3 --- Sprint Planner.}\label{fig:burndown-s3}
\end{figure}

\begin{figure}[H]\centering
\includegraphics[width=0.85\linewidth]{burndown-s4}
\caption{Burndown du sprint 4 --- Epic Grooming et routage multi-modèles.}\label{fig:burndown-s4}
\end{figure}

\begin{figure}[H]\centering
\includegraphics[width=0.85\linewidth]{burndown-s5}
\caption{Burndown du sprint 5 --- UX et accessibilité.}\label{fig:burndown-s5}
\end{figure}

\begin{figure}[H]\centering
\includegraphics[width=0.85\linewidth]{burndown-s6}
\caption{Burndown du sprint 6 --- Déploiement et démonstration.}\label{fig:burndown-s6}
\end{figure}
